\documentclass[11pt]{article}

\usepackage[T1]{fontenc}
\usepackage{lmodern}
\usepackage[margin=1.15in]{geometry}
\usepackage{amsmath,amssymb,amsthm}
\usepackage{mathtools}
\usepackage{xcolor}
\usepackage{framed}
\usepackage{array}
\usepackage{booktabs}
\usepackage[section]{placeins}
\usepackage{enumitem}
\usepackage[protrusion=true,expansion=false]{microtype}
\usepackage{xurl}
\usepackage[colorlinks=true,linkcolor=blue!50!black,citecolor=blue!50!black,urlcolor=blue!50!black]{hyperref}
\hypersetup{
  pdftitle={An upper bound for the de Bruijn--Newman constant from a zero-free half-plane},
  pdfauthor={Jude Gomila},
  pdfsubject={Computer-assisted upper bound for the de Bruijn--Newman constant (draft, not peer reviewed)},
  pdfkeywords={de Bruijn--Newman constant, Riemann xi function, heat flow, zero-free region, barrier method, Lean}
}

% ---------------------------------------------------------------
% Theorem environments
% ---------------------------------------------------------------
\newtheorem{theorem}{Theorem}[section]
\newtheorem{lemma}[theorem]{Lemma}
\newtheorem{proposition}[theorem]{Proposition}
\newtheorem{corollary}[theorem]{Corollary}

\theoremstyle{definition}
\newtheorem{definition}[theorem]{Definition}
\newtheorem{remark}[theorem]{Remark}

\newcommand{\LB}{\Lambda}
\newcommand{\R}{\mathbb{R}}
\newcommand{\C}{\mathbb{C}}
\newcommand{\capo}{\mathrm{cap}}
\newcommand{\qe}{q_{\mathrm e}}
\DeclareMathOperator{\Rea}{Re}
\DeclareMathOperator{\Ima}{Im}

\emergencystretch=2em

\title{An upper bound for the de Bruijn--Newman constant\\
from a zero-free half-plane}
\author{Jude Gomila}
\date{Draft, 9 October 2026}

\begin{document}
\maketitle

\begin{abstract}
\noindent
Let $\LB$ be the de Bruijn--Newman constant. Assume that $\zeta(s)\ne0$ for
$\Rea s>7/8$, that the analytic statements of D.~J.~Raymond's decreasing-barrier
construction hold, and that the Riemann hypothesis holds for zeros of height at most
$3000175332800$ (Platt and Trudgian). Then
\[
\LB\le\frac{118505710922701}{10^{15}}=0.118505710922701 .
\]
If in addition the floor and kernel certificates of a re-tuned run produced for this
paper hold, then
\[
\LB\le\frac{116091488672439}{10^{15}}=0.116091488672439 .
\]
The half-plane enters Raymond's chain only through the starting value of the barrier
and its strict-entry check, which become $q_0>(2\theta-1)^2=9/16$ in place of $q_0>1$.
The half-plane $\Rea s>7/8$ is a claim of the openai/math repository whose Lean proof
was rebuilt locally; it has not been peer reviewed, and Raymond's construction has
not been refereed. The paper also gives the value of the same argument as a function of the
zero-free abscissa $\theta$, records what limits the abscissa $7/8$, and reports a
Lean~4 check of the deduction from named hypotheses in which every step of both runs is
replayed by the kernel. The accepted peer-reviewed bound remains $\LB\le0.2$.
\end{abstract}

\begin{center}
\begin{minipage}{0.88\textwidth}
\begin{framed}
\small
\raggedright
\noindent\textbf{Status of this draft.}
A computer-assisted upper bound. It depends on two recent results that have not been
refereed: the zero-free half-plane $\Rea s>7/8$ (openai/math result family~003, whose
Lean proof was rebuilt here with standard axioms only) and D.~J.~Raymond's
decreasing-barrier construction (Zenodo \href{https://doi.org/10.5281/zenodo.22867547}{10.5281/zenodo.22867547}).
The value $0.116091$ also uses this project's internal floor and kernel certificates;
without them the same argument gives $0.118506$. Theorem~\ref{thm:main} states these
inputs as hypotheses (H1)--(H4) (Section~\ref{sec:main}). Not submitted, not published,
not yet peer reviewed. The accepted peer-reviewed bound remains $\LB\le0.2$
(Platt--Trudgian \cite{PlattTrudgian21}). The Riemann hypothesis is open.
\end{framed}
\end{minipage}
\end{center}

% ===============================================================
\section{Introduction}
% ===============================================================

\subsection{The de Bruijn--Newman constant}

Let
\[
\Phi(u)=\sum_{n=1}^{\infty}\bigl(2\pi^2n^4e^{9u}-3\pi n^2e^{5u}\bigr)\exp\bigl(-\pi n^2e^{4u}\bigr),
\qquad
H_t(z)=\int_0^\infty e^{tu^2}\Phi(u)\cos(zu)\,du
\quad(t\in\R,\ z\in\C).
\]
With $\xi(s)=\tfrac12s(s-1)\pi^{-s/2}\Gamma(s/2)\zeta(s)$ one has
$H_0(z)=\tfrac18\,\xi\bigl(\tfrac12+\tfrac{iz}2\bigr)$, so the zeros of $H_0$ are the
points $z=2\gamma-i(2\beta-1)$ for the nontrivial zeros $\rho=\beta+i\gamma$ of $\zeta$.
By the work of de~Bruijn \cite{deBruijn50} and Newman \cite{Newman76} there is a finite
constant $\LB\le\frac12$, the \emph{de Bruijn--Newman constant}, such that all zeros of
$H_t$ are real exactly when $t\ge\LB$. The Riemann hypothesis (RH) is equivalent to
$\LB\le0$. Newman conjectured $\LB\ge0$; Rodgers and Tao proved it \cite{RodgersTao20}.
An upper bound on $\LB$ therefore measures how far RH can fail in this normalisation;
it does not prove RH.

\subsection{Upper bounds}\label{sec:history}

De~Bruijn proved $\LB\le\frac12$ \cite{deBruijn50}, and Ki, Kim and Lee proved
$\LB<\frac12$ \cite{KiKimLee09}. The Polymath project developed effective
approximations of $H_t$ and a barrier criterion, and proved $\LB\le0.22$
\cite{Polymath19}. Platt and Trudgian verified RH up to height $3000175332800$ and
observed that a row of the Polymath parameter table then gives $\LB\le0.2$
\cite[Theorem~1 and Corollary~2]{PlattTrudgian21}. This is the current peer-reviewed
upper bound. Histories of the earlier bounds, including the lower bounds that preceded
\cite{RodgersTao20}, are given in \cite{RodgersTao20} and \cite{Polymath19}.

In 2026 several preprints have claimed smaller upper bounds without a zero-free half-plane:
$0.1875$ \cite{Mosaic26}, $0.1787854$ \cite{Gomila26}, $0.163297$ \cite{Raymond26a},
$0.157656619\ldots$ \cite{Gordon26} and $0.149396$ \cite{Raymond26}. None of these has
been peer reviewed. The present paper uses the method of \cite{Raymond26} and no
result of the others.

\subsection{Main result}\label{sec:main}

Fix $X=6000000185827$. The hypotheses are the following; Sections~\ref{sec:raymond}
and~\ref{sec:computation} give the precise statements.

\begin{enumerate}[leftmargin=3.2em]
\item[\textbf{(H1)}] $\zeta(s)\ne0$ for every $s$ with $\Rea s>7/8$.
\item[\textbf{(H2)}] The statements of Raymond's construction \cite{Raymond26} listed in
Section~\ref{sec:raymond}: the head localisation (W), the probe identity and zero
dynamics used in Lemma~\ref{lem:raymond}, the floor inequalities (J) for the entries of
his floor tables used by the run of Section~\ref{sec:run1}, and the kernel inequality (K)
for his catalogue kernels used by that run.
\item[\textbf{(H3)}] Every zero $\beta+i\gamma$ of $\zeta$ with $0<\gamma\le3000175332800$
has $\beta=\frac12$ \cite[Theorem~1]{PlattTrudgian21}.
\item[\textbf{(H4)}] The floor inequalities (J) for the $468$ floor entries and the kernel
inequality (K) for the $59$ kernel rows used by the run of Section~\ref{sec:run2}.
\end{enumerate}

\begin{theorem}\label{thm:main}
\begin{enumerate}[label=\textup{(\roman*)},leftmargin=2.2em]
\item Assume \textup{(H1)}, \textup{(H2)} and \textup{(H3)}. Then
\[
\LB\le\frac{118505710922701}{10^{15}}=0.118505710922701 .
\]
\item Assume \textup{(H1)}--\textup{(H4)}. Then
\[
\LB\le\frac{116091488672439}{10^{15}}=0.116091488672439 .
\]
\end{enumerate}
\end{theorem}

Part (ii) uses from (H2) only (W) and the identities of Lemma~\ref{lem:raymond}, not
Raymond's floor tables; its floors and kernels are those of (H4).

\paragraph{Status of the hypotheses.}
(H1) is the claim of result family~003 of the openai/math repository
\cite{OpenAI26QRH}, whose manuscripts the repository describes as produced by an
internal OpenAI model. Its Lean proof builds locally with only the standard axioms and
its statement matches the repository's comparator challenge
(Section~\ref{sec:lean}); it has not been peer reviewed. (H2) is part of Raymond's
unrefereed preprint \cite{Raymond26}; it has been read only inside this project, in two
internal review passes \cite{IntReviewA,IntReviewB}. (H3) is peer reviewed. (H4)
consists of this project's certificates, most of them single implementations, with one
internal referee pass.

\paragraph{Dependency.}
Remove (H1) and the same chain gives Raymond's value $0.149395073827688$
\cite{Raymond26}. Remove (H2) or (H3) and the argument gives nothing below the accepted
$0.2$: (H1) alone gives $\LB\le(2\cdot\frac78-1)^2/2=\frac9{32}=0.28125$
(Lemma~\ref{lem:subst}). Remove (H4) and (ii) reduces to (i). Each of (H1), (H2) and (H3)
is needed for (i), and (H4) only for (ii).

\subsection{Raymond's method, and the one change made here}

The method is Raymond's \cite{Raymond26}. His construction controls the squared height
$q=y^2$ of the highest zero $x+iy$ of $H_t$ to the right of $X$ by a decreasing barrier
\[
\beta'(t)=-2-c\,\beta(t),
\]
where the rates $c>0$ come from certified lower bounds for $\partial_y\log|H_t(x+iy)|$
at two translated probe points, combined with exact kernel inequalities. A localisation
wall, built on Platt and Trudgian's verification of RH, keeps every nonreal zero to the
right of $X$. The barrier starts at time $0$ above the classical strip, $q_0>1$; the bound
on $\LB$ is the time at which the barrier reaches $10^{-8}$, plus $5\cdot10^{-9}$. His
floor estimates, kernel catalogue, choice rule and exact comparison arithmetic give
$\LB\le18674384228461/(125\cdot10^{12})=0.149395073827688<0.149396$.

This paper makes one change: if $\zeta(s)\ne0$ for $\Rea s>\theta$, the barrier may start
at $q_0>(2\theta-1)^2$. For $\theta=7/8$ the start drops from just above $1$ to just above
$9/16$, and the barrier reaches $10^{-8}$ about $0.031$ earlier. Lemma~\ref{lem:subst}
proves the time-zero bound and Theorem~\ref{thm:barrier} proves that the start and its
strict-entry check are the only places where the zero locations at time $0$ enter the
chain. The computation of Section~\ref{sec:run1} then reruns Raymond's chain unchanged
apart from the start. The run of Section~\ref{sec:run2} uses, in addition, floors and
kernels produced in this project, re-tuned along the lower barrier path.

The substitution applies to any decreasing-barrier chain of this type. Raymond's own
chain with the $7/8$ start gives part~(i) with no input from this project other than Lemma~\ref{lem:subst} and the written completion of the first-contact
argument in Theorem~\ref{thm:barrier}.

\subsection{Further results}

Section~\ref{sec:theta} gives the value of the argument at the Platt--Trudgian height as a
function of a hypothetical zero-free abscissa $\theta$: near $\theta=7/8$, each $0.01$
of $\theta$ is worth about $0.003$ of $\LB$. Section~\ref{sec:limits} records the
internal analyses of what limits the abscissa $7/8$ in \cite{OpenAI26QRH}; none of them
improves it. Section~\ref{sec:lean} describes a Lean~4 development in which
Theorem~\ref{thm:main} is proved from named hypotheses, with every arithmetic step of
both runs replayed by the kernel.

\paragraph{Use of generative AI.}
The half-plane (H1) was produced by an internal OpenAI model, as stated in
\cite{OpenAI26QRH}. The computations, the internal reviews, the Lean developments and
the text of this draft were produced with a multi-model approach, using several
generative AI systems in human-directed cycles of exploration, programming, checking
and revision.
Jude Gomila directed the work and is responsible for the claims and the text.

% ===============================================================
\section{Zero dynamics under the heat flow}\label{sec:dynamics}
% ===============================================================

The function $H_t$ is even, real on the real axis and of order~$1$, and $H_t(iy)>0$ for
real $y$ \cite[Section~3]{Polymath19}. Its zeros are therefore symmetric under
$z\mapsto-z$ and $z\mapsto\bar z$, and none lies on the imaginary axis. Write
$z=x+iy$. The following results of \cite{Polymath19} are used.

\begin{proposition}[{\cite[Proposition~3.1]{Polymath19}}]\label{prop:dyn}
Let $t_0\in\R$.
\begin{enumerate}[label=\textup{(\roman*)},leftmargin=2.2em]
\item If $z_j(t_0)$ is a simple zero of $H_{t_0}$, it extends to a smooth family of simple
zeros $z_j(t)$ for $t$ near $t_0$, with
$\frac{d}{dt}z_j(t_0)=2\sideset{}{'}\sum_{k\ne j}\frac1{z_j(t_0)-z_k(t_0)}$, the terms
$k$ and $-k$ being summed together.
\item If $z_j(t_0)$ is a zero of order $m\ge2$, then for $t$ close to $t_0$ the $m$ zeros
of $H_t$ near it are
$z_j(t_0)+\sqrt2\,(t-t_0)^{1/2}\lambda_l+O(|t-t_0|)$, $l=1,\dots,m$, where
$\lambda_1<\dots<\lambda_m$ are the roots of the Hermite polynomial $\mathrm{He}_m$.
\end{enumerate}
\end{proposition}

\begin{theorem}[{de~Bruijn \cite[Theorem~13]{deBruijn50}; \cite[Theorem~3.2]{Polymath19}}]\label{thm:debruijn}
Let $t_0\in\R$ and $y_0>0$, and suppose $H_{t_0}$ has no zeros with $y>y_0$. Then for
every $t>t_0$, $H_t$ has no zeros with $y>\max(y_0^2-2(t-t_0),0)^{1/2}$. In particular
$\LB\le t_0+\frac12y_0^2$.
\end{theorem}

\begin{theorem}[{\cite[Theorem~1.5(i)]{Polymath19}}]\label{thm:confine}
There is an absolute constant $C>0$ such that for $0<t\le\frac12$ every zero $x+iy$ of
$H_t$ with $x\ge\exp(C/t)$ is real.
\end{theorem}

% ===============================================================
\section{Raymond's decreasing-barrier method}\label{sec:raymond}
% ===============================================================

This section states the parts of \cite{Raymond26} that are used. Section references
``AF'' are to the companion \emph{Analytic foundations for the time-zero upper-bound
construction} in the same deposit; ``main'' refers to the paper \emph{Correlated
translated probes}. The fixed data are
\[
X=6000000185827,\qquad N_0=690988,\qquad \qe=10^{-8},
\]
the $608$ closed time cells $C_i=[i/4000,(i+1)/4000]$, $0\le i\le607$, covering
$[0,19/125]$, and a probe height $h_i$ for each cell (Raymond's plan uses rationals in
$[0.69,2.18]$). A zero $z$ of $H_t$ is a \emph{tail zero} if $|\Rea z|>X$.

\subsection{The head wall}

\begin{enumerate}[leftmargin=3em]
\item[\textbf{(W)}] For $0\le t\le41/250$, every zero $z$ of $H_t$ with $|\Rea z|\le X$ is
real.
\end{enumerate}
This is AF~\S\S2--3. Platt and Trudgian's verification (H3) gives reality of the zeros
of $H_0$ with $|\Rea z|<X$; the ordinate used is $(X+1)/2=3000000092914$, which lies
$175239886$ below $3000175332800$. An interval certificate proves
$|H_t(X+iy)/B_t(X+iy)|>3/4$ for $0\le t\le41/250$, $0\le y\le1$, where $B_t$ is the
normalising factor of \cite{Polymath19}; with the classical strip $|y|\le1$ this closes
the head cell $[-X,X]\times[-1,1]$, and a counting argument with
Proposition~\ref{prop:dyn} shows that no nonreal zero appears inside it. This is the only
use of (H3).

\subsection{Probe floors}

For $h>0$ put $S(x,h,t)=\partial_y\log|H_t(x+iy)|\big|_{y=h}$ and
\[
K_h(d,b)=\frac{h-b}{d^2+(h-b)^2}+\frac{h+b}{d^2+(h+b)^2}\qquad(d\in\R,\ 0\le b).
\]
Since $H_t$ is even of order~$1$ with $H_t(0)>0$, its Hadamard product gives
\begin{equation}\label{eq:probe}
S(x,h,t)=\sum_{\text{pairs}}K_h(x-x_j,b_j),
\end{equation}
summed over the conjugate pairs $x_j\pm ib_j$ of zeros of $H_t$, with real zeros counted
with weight $\frac12$; the sum converges absolutely since $K_h(d,b)\ll h/d^2$.

\begin{enumerate}[leftmargin=3em]
\item[\textbf{(J)}] A \emph{floor entry} $(i,r,h,\ell)$ is valid if for every $t\in C_i$ and
every $x\ge X$,
\[
\tfrac12\bigl[S(x+rh,h,t)+S(x-rh,h,t)\bigr]>\ell .
\]
\end{enumerate}
Raymond's floors come from interval computations built on his one-point estimates (AF
\S\S4--9) and his two-probe joint floor (main, Theorem~4). They are stated for every
$x\ge X$, including $x\to\infty$, with six frequency ranges per cell, the last one
infinite.

\subsection{Kernels and the self term}

\begin{enumerate}[leftmargin=3em]
\item[\textbf{(K)}] A \emph{kernel row} $(r,\lambda,\rho)$ is valid if for every $h>0$, every
real $d$, every $0<y\le\rho h$ and every $0\le b\le y$ with $(d,b)\ne(0,y)$,
\[
2K_y(d,b)\ \ge\ \lambda\,\frac yh\cdot\frac{K_h(d+rh,b)+K_h(d-rh,b)}2 .
\]
\end{enumerate}
This is AF~(10.1). Raymond proves it for a catalogue of $13$ rows by exact Bernstein
expansions (below, ``catalogue kernel~$j$'' is the row with index $j$, counted from $0$, in
the order of his \texttt{data/KERNELS.json}; AF does not number the rows); the inequality is invariant under the scaling $(d,b,y,h)\mapsto(cd,cb,cy,ch)$.
For a zero at height $y$ write $q=y^2$ and $v=q/h^2$. The zero's own conjugate pair
contributes $K_h(\pm rh,y)=s_r(v)/h$ to each probe, where
\[
s_r(v)=\frac{2(1+r^2-v)}{(1+r^2-v)^2+4r^2v}.
\]
A \emph{cap} for $(r,v)$ is a number $\capo\ge\sup_{0\le v'\le v}s_r(v')$, and the
\emph{rate} of a floor entry with ratio $r$ and a kernel row $(r,\lambda,\rho)$ is
\[
c=\frac{2\lambda}h\Bigl(\ell-\frac{\capo}h\Bigr).
\]

\subsection{The barrier inequality}

\begin{lemma}[{AF~(11.1)}]\label{lem:raymond}
Let $t>0$ lie in a cell $C_i$, and let $z^*=x^*+iy^*$ be a simple zero of $H_t$ with
$x^*>X$, $y^*>0$, and $y^*\ge|\Ima w|$ for every zero $w$ of $H_t$. Let $(i,r,h,\ell)$ be a
valid floor entry and $(r,\lambda,\rho)$ a valid kernel row with $h=h_i$, and suppose
$v=y^{*2}/h^2\le\rho^2$ and $\capo\ge\sup_{[0,v]}s_r$. Then $q(t)=(\Ima z^*(t))^2$
satisfies
\[
q'(t)<-2-c\,q(t),\qquad c=\frac{2\lambda}h\Bigl(\ell-\frac{\capo}h\Bigr).
\]
\end{lemma}

\begin{proof}
Grouping the sum in Proposition~\ref{prop:dyn}(i) into conjugate pairs (the imaginary
parts converge absolutely), the conjugate $\bar z^*$ contributes $-1/y^*$ to
$\frac{d}{dt}\Ima z^*$, and every other pair $x_j\pm ib_j$ contributes
$-2K_{y^*}(x^*-x_j,b_j)$, real zeros with weight $\frac12$. Hence
$q'=-2-4y^*\sum'K_{y^*}(d_j,b_j)$ with $d_j=x^*-x_j$, the sum over the other pairs. Every
other zero has $b_j\le y^*$, and $(d_j,b_j)\ne(0,y^*)$ because $z^*$ is simple. Applying (K)
to each term and using \eqref{eq:probe} at $x^*\pm rh$, where the own pair contributes
$s_r(v)/h$ to each probe,
\[
4y^*\sideset{}{'}\sum K_{y^*}(d_j,b_j)\ \ge\ \frac{2\lambda q}h\Bigl(\tfrac12\bigl[S(x^*+rh,h,t)+S(x^*-rh,h,t)\bigr]-\frac{s_r(v)}h\Bigr)
\ >\ \frac{2\lambda q}h\Bigl(\ell-\frac{\capo}h\Bigr)
\]
by (J) at $x=x^*\ge X$ and $s_r(v)\le\capo$. Hence $q'<-2-cq$.
\end{proof}

\subsection{Discretisation and the terminal step}

Each cell is divided into $128$ substeps of length $\delta=1/512000$; write $t_k=k\delta$
and $i(k)=\lfloor k/128\rfloor$. A run is a record $(q_k,c_k)_{0\le k\le K}$ with a mode
(ratio $r_k$, kernel row $(r_k,\lambda_k,\rho_k)$, floor $\ell_k$ for cell $i(k)$, cap) on
each substep. The exact solution of $\beta'=-2-c\beta$ over one substep is
$q\mapsto(q+2/c)e^{-c\delta}-2/c$. The stored values are upper enclosures: $e^{-x}$ is
replaced by its degree-$18$ Taylor polynomial (an upper bound for $0\le x<1$) and the
result is rounded up to a multiple of $2^{-96}$. At the terminal substep the remaining
time $\tau=\log\bigl(1+c(q-\qe)/(2+c\qe)\bigr)/c$ is bounded above with the alternating
Taylor sum of $\log(1+z)$ through degree $11$, and
\[
U=\bigl\lceil t_K+\tau+\qe/2\bigr\rceil_{10^{-15}} .
\]
At each substep the choice rule takes, among the admissible modes, one with the largest
certified rate; any admissible choice gives a valid barrier.

With the classical start $q_0=(1000001/1000000)^2$, Raymond's run gives
\cite[Theorem~1]{Raymond26}
\[
\LB\le\frac{18674384228461}{125000000000000}=0.149395073827688<0.149396 .
\]

% ===============================================================
\section{The substitution lemma}\label{sec:subst}
% ===============================================================

\begin{lemma}\label{lem:subst}
Let $\frac12<\theta<1$ and assume $\zeta(s)\ne0$ for $\Rea s>\theta$. Then every zero $z$ of
$H_0$ satisfies $|\Ima z|\le2\theta-1$, and for every $t\ge0$ every zero $z$ of $H_t$
satisfies
\[
(\Ima z)^2\le\max\bigl((2\theta-1)^2-2t,\ 0\bigr).
\]
In particular $\LB\le(2\theta-1)^2/2$.
\end{lemma}

\begin{proof}
Write $s=\frac12+\frac{iz}2$, so that $\Rea s=(1-\Ima z)/2$. For $\Rea s>\theta\ge\frac12$
the factors $s$, $\pi^{-s/2}$ and $\Gamma(s/2)$ of $\xi$ do not vanish; $\xi(1)=\frac12$;
and for $s\ne1$, $s-1\ne0$ and $\zeta(s)\ne0$ by hypothesis. So $\xi$ has no zero with
$\Rea s>\theta$, and by $\xi(s)=\xi(1-s)$ none with $\Rea s<1-\theta$. Every zero of $H_0$
therefore has $1-\theta\le\Rea s\le\theta$, that is,
$|\Ima z|=|1-2\Rea s|\le2\theta-1$. Theorem~\ref{thm:debruijn} with $t_0=0$ and
$y_0=2\theta-1$ gives the bound for $\Ima z\ge0$, and the symmetry $z\mapsto\bar z$ gives
it for $\Ima z<0$. At $t=(2\theta-1)^2/2$ all zeros are real.
\end{proof}

For $\theta=7/8$ this gives $(\Ima z)^2\le\max(\frac9{16}-2t,0)$ and $\LB\le\frac9{32}$. The bound is
closed: zeros with $\beta=7/8$ are not excluded, so $\Ima z=\frac34$ can occur at $t=0$.
Lemma~\ref{lem:subst} is proved in Lean (Section~\ref{sec:lean}).

The next theorem is Raymond's first-contact argument (AF \S11; main \S7) with the start
condition $q_0>1$ replaced by $q_0>(2\theta-1)^2$, written out in full. It combines
Lemma~S($\theta$) and Lemma~FB($\theta$) of the internal review \cite[\S\S5,
8]{IntReviewB}; the second internal review confirms the start substitution
independently \cite[\S7]{IntReviewA}.

\begin{theorem}\label{thm:barrier}
Let $\frac12<\theta<1$. Assume $\zeta(s)\ne0$ for $\Rea s>\theta$, assume \textup{(W)}, and
let $(q_k,c_k)_{0\le k\le K}$ be a run whose floor entries satisfy \textup{(J)} and whose
kernel rows satisfy \textup{(K)}. Write $h_k=h_{i(k)}$ and $v_k=q_k/h_k^2$, and suppose:
\begin{enumerate}[label=\textup{(C\arabic*)},leftmargin=3em]
\item $v_k\le\rho_k^2$ and $r_kh_k<2$;
\item $c_k=\frac{2\lambda_k}{h_k}\bigl(\ell_k-\capo_k/h_k\bigr)>0$ with
$\capo_k\ge\sup_{0\le v\le v_k}s_{r_k}(v)$;
\item $0<c_k\delta<1$ and $q_{k+1}\ge(q_k+2/c_k)e^{-c_k\delta}-2/c_k$ for $k<K$;
\item \textup{(S1)} $q_0-(2+c_0q_0)\cdot10^{-8}>(2\theta-1)^2$;
\item $(q_K+2/c_K)e^{-c_K\delta}-2/c_K\le\qe$, and
$\tau=\log\bigl(1+c_K(q_K-\qe)/(2+c_K\qe)\bigr)/c_K\le\delta$;
\item $T^*=t_K+\tau$ satisfies $T^*+\qe/2<41/250$.
\end{enumerate}
Then $\LB\le T^*+\qe/2$.
\end{theorem}

\begin{proof}
\emph{Step 1 (barrier).} Let $\beta$ be the continuous function with $\beta(0)=q_0$ and
$\beta'=-2-c_k\beta$ on $[t_k,t_{k+1}]$. The one-step map is increasing in $q$, so by (C3)
and induction $\beta(t_k)\le q_k$; since $\beta$ decreases on each substep,
$\beta(t)\le q_k$ for $t\in[t_k,t_{k+1}]$. Conditions (C1)--(C2) are upper bounds on the
height (admissibility $v\le\rho^2$; the cap is a supremum over $[0,v]$), so they hold at
every height at most $\beta(t)$ on substep $k$. By (C5) and monotonicity of the one-step
map, $\beta$ reaches $\qe$ at a time $T_\beta\le T^*$.

\emph{Step 2 (start, $0\le t\le10^{-8}$).} On the first substep $\beta$ is convex and
decreasing with $\beta'(0)=-2-c_0q_0$, so $\beta(t)\ge q_0-(2+c_0q_0)t$. Since
$10^{-8}<\delta$, Lemma~\ref{lem:subst} and (S1) give, for $0\le t\le10^{-8}$ and every zero,
\[
(\Ima z)^2\le(2\theta-1)^2-2t\le(2\theta-1)^2<q_0-(2+c_0q_0)\cdot10^{-8}\le\beta(t).
\]
No zero meets the barrier before time $10^{-8}$. This is the only step that uses the
zero locations at time $0$, and the only use of the half-plane.

\emph{Step 3 (no contact on $[10^{-8},T_\beta]$).} Let $E$ be the set of
$t\in[10^{-8},T_\beta]$ at which some tail zero satisfies $(\Ima z)^2\ge\beta(t)$, and
suppose $E\ne\emptyset$. For $t\in[10^{-8},\frac12]$, Theorem~\ref{thm:confine} and
$|\Ima z|\le1$ (Theorem~\ref{thm:debruijn} with $t_0=0$, $y_0=1$) place every nonreal zero
in a fixed compact set. By Hurwitz's theorem, the continuity of $\beta$ and (W) (a limit
of tail zeros that is nonreal is a tail zero), $E$ is closed; let $t^*=\min E$. By
Step~2 and closedness, $t^*>10^{-8}$. Every tail zero at time $t^*$ is a limit of zeros
at earlier times, at which $(\Ima z)^2<\beta$; hence some tail zero $z^*$ satisfies
$(\Ima z^*)^2=\beta(t^*)$ and no zero is higher. By (W) and the symmetries we may take
$x^*>X$ and $y^*>0$. Let $k$ be the substep with $t^*\in(t_k,t_{k+1}]$ (the left substep if
$t^*$ is a knot).

If $z^*$ is a repeated zero, Proposition~\ref{prop:dyn}(ii) with $t<t^*$ gives zeros
$z^*+i\sqrt{2(t^*-t)}\,\lambda_l+O(t^*-t)$ near $z^*$; since $\lambda_m>0$, one of them has
$\Ima z-y^*\gg(t^*-t)^{1/2}$, while $\sqrt\beta$ is Lipschitz near $t^*$ (as $\beta\ge\qe>0$).
That zero violates $(\Ima z)^2<\beta(t)$ for $t<t^*$ close to $t^*$, contrary to the
minimality of $t^*$.

If $z^*$ is simple, Step~1 gives $v^*=\beta(t^*)/h_k^2\le v_k\le\rho_k^2$ and
$\capo_k\ge s_{r_k}(v^*)$, so Lemma~\ref{lem:raymond} applies (the cell of substep $k$ is
closed and contains $t^*$) and gives $q'(t^*)<-2-c_k\beta(t^*)=\beta'(t^*-)$. But
$q(t)<\beta(t)$ for $t<t^*$ and $q(t^*)=\beta(t^*)$ force $q'(t^*)\ge\beta'(t^*-)$. This is a
contradiction, so $E=\emptyset$.

\emph{Step 4 (terminal).} By Step~3 every tail zero at time $T_\beta$ satisfies
$(\Ima z)^2<\qe$, and by (W) and $T_\beta<41/250$ every other zero is real.
Theorem~\ref{thm:debruijn} with $t_0=T_\beta$ and $y_0=\sqrt{\qe}$ gives
$\LB\le T_\beta+\qe/2\le T^*+\qe/2$.
\end{proof}

\begin{remark}[Where the half-plane enters]\label{rem:where}
The half-plane is used once, in Step~2. The wall (W) uses the classical strip and (H3);
(H1) bounds the head zeros by $|\Ima z|\le\frac34$ but does not make them real, so it does
not replace (H3). The floors (J) and kernels (K) read no zero locations. The confinement
after time $10^{-8}$ uses Theorem~\ref{thm:confine} and $|\Ima z|\le1$. The terminal
allowance $\qe/2$ comes from Theorem~\ref{thm:debruijn} at $T_\beta$ and does not depend on
$q_0$. For $\theta<1$ the barrier starts below $1$ and falls faster than the classical
envelope $1-2t$, so on the whole run it lies below $1-2t$ and the classical strip never
keeps the zeros under it. After time $10^{-8}$ the invariant is carried by the
first-contact argument alone.
\end{remark}

\begin{remark}[Windows]
Theorem~\ref{thm:barrier} uses a barrier that is continuous across substeps. A floor valid
only on a window $q\in[q_{\mathrm{lo}},q_{\mathrm{hi}}]$ would need the stored, left-continuous
form of the barrier \cite[\S\S4.2, 8]{IntReviewB}. Neither run uses such a floor.
\end{remark}

\begin{remark}[The fixed-time criterion]
In the Polymath criterion \cite[Theorem~1.2]{Polymath19} the half-plane lowers the tops of
the regions in hypotheses (ii) and (iii) to $\sqrt{(2\theta-1)^2-2t}$, but the conclusion
$\LB\le t_0+y_0^2/2$ and its binding height $y_0$ are unchanged. On this reading
(a method-level analysis, not a theorem of this paper), a zero-free half-plane does not move
a bound of the Polymath type directly; it moves a decreasing-barrier bound through its start.
\end{remark}

\begin{remark}[Classical zero-free regions]
Zero-free regions near $\Rea s=1$ do not give a start below $1$. For example, the region
$\sigma\ge1-1/(5.558691\log|t|)$ of Mossinghoff, Trudgian and Yang \cite{MTY24} gives
$\beta<0.99373824$ at $\gamma=3.0001753\cdot10^{12}$, so $(2\beta-1)^2<0.9751098$ there, and
its width tends to $0$ as $\gamma\to\infty$; it gives no uniform bound $q_0<1$. Given (H3),
a half-plane valid only above the verified height is the same statement as a global one.
\end{remark}

% ===============================================================
\section{The computation}\label{sec:computation}
% ===============================================================

Both runs use exact rational arithmetic for the barrier and its checks. The floor tables
are interval certificates computed with Arb \cite{JohanssonArb17}, through python-flint
\cite{FLINT} and, for the structured floors of Section~\ref{sec:run2}, also through a C
program with complex balls.
Numbers below are from the run records listed in Appendix~\ref{app:repro}.

\subsection{Value (i): Raymond's chain with the 7/8 start}\label{sec:run1}

The run uses Raymond's floor tables at the Platt--Trudgian height (regenerated
bit-identically from his code), his plan heights and his kernel catalogue, with
\[
q_0=\frac{281251}{500000}=\frac9{16}+2\cdot10^{-6}.
\]
On the first substep the chosen mode is catalogue kernel~4 ($r=23/40$, $\lambda=199/100$,
$\rho=7/20$) at $h=109/50$, with $v_0=q_0/h^2=0.118362\le\rho^2=0.1225$, cell-0 floor
$\ell=627382516099/10^{11}$ and global cap $1503067131/10^9$, so
\[
c_0=\frac{12113002361358409}{1188100000000000}\approx10.195272,
\]
\[
q_0-(2+c_0q_0)\cdot10^{-8}-\frac9{16}=\frac{114215105972865586110341}{59405\cdot10^{24}}\approx1.922651\cdot10^{-6}.
\]
The run has $60675$ substeps. It uses only catalogue kernels~4 ($51614$ substeps) and~5
($r=3/5$, $\lambda=199/100$, $\rho=1/5$; $9061$ substeps), and never Raymond's sharp
continuum law (AF Part~II). The tightest admissibility margin on the path is
$\rho^2-v=3.793\cdot10^{-6}$. The terminal substep is $k=60674$ (index $2$ within cell
$474$), $t_K=30337/256000$, with $z=3.4187\cdot10^{-5}$ and $\tau/\delta=0.9214$. The result is
\[
U=\frac{118505710922701}{10^{15}}=0.118505710922701 .
\]
Raymond's alternate logarithm bound $\log(1+z)\le z-z^2/(2(1+z))$ gives
$59252855461351/(5\cdot10^{14})=0.118505710922702$; both are valid upper bounds and the
smaller is used. Three separately written implementations reproduce the run bit for bit: the project's
replay of Raymond's consumer, an implementation by the internal referee that shares no
code with it (with an exact-rational lane and a $320$-bit Arb lane), and the replay written
for the internal review \cite[\S7]{IntReviewB}; the Lean kernel recomputes the run as well
(Section~\ref{sec:lean}).

\subsection{Value (ii): a re-tuned run with additional certificates}\label{sec:run2}

The second run replaces Raymond's floor tables by a larger menu of certified floors and
kernels produced in this project (called \emph{route~P} in the project's records), and
re-tunes the menu along the lower barrier path that the $7/8$ start produces. The menu
consists of:
\begin{enumerate}[label=(M\arabic*),leftmargin=2.6em]
\item joint floors in which Raymond's two-probe surplus is retained for further label ranges
(an extension of main Theorem~4 to arbitrary shift intervals), computed with his producer
at re-planned probe heights between $0.86$ and $1.07$ times his plan heights: $289$ floor
entries, used on $36923$ substeps;
\item joint floors from the project's Theorem~MP (a joint floor for a weighted probe measure,
which for a single pair of probes is main Lemma~3), at Raymond's plan heights: $90$ entries,
$11212$ substeps;
\item structured floors in which the primes up to $47$ are treated exactly rather than
bounded term by term, at $0.94$ and $1$ times the plan heights: $89$ entries, $11304$
substeps;
\item $79$ new two-point kernels with $r=3/5$ and $\lambda<2$ on the grid
$\rho\in\{81/400,\dots,171/400\}$, certified by exact Bernstein expansions at subdivision
depth $22$, of which $50$ are used; together with $6$ earlier project kernel rows and
Raymond's catalogue kernels $4$, $8$, $10$, the run uses $59$ kernel rows.
\end{enumerate}
The run starts at
\[
q_0=\frac{5625001}{10^7}=\frac9{16}+10^{-7},\qquad
c_0=\frac{41336831094003950997}{3755300000000000000}\approx11.0076,
\]
\[
q_0-(2+c_0q_0)\cdot10^{-8}-\frac9{16}=\frac{67904283759396681637924003}{37553\cdot10^{29}}\approx1.8082\cdot10^{-8}.
\]
It has $59439$ substeps, ends with substep $k=59438$ (index $46$ within cell $464$),
$t_K=29719/256000$, uses probe heights in
$[91/100,47/25]$ and no window floor, and gives
\[
U=\frac{116091488672439}{10^{15}}=0.116091488672439,
\]
equal to the value obtained with the alternate logarithm bound.

Table~\ref{tab:ladder} lists the re-tuning steps, each an exact run. The final value is
computed by its own run, not by adding the steps. About $74\%$ of the gain comes from the
lowered heights (rungs E3 and E4), $24.5\%$ from the existing height tables of rung E2 and
$1.2\%$ from the new kernels. Lowering the heights further is limited by Raymond's producer:
its moment cap requires $C<17/5$ for the exceptional labels in the head $n\le1024$, which
fails at $0.84$ times the plan in every scanned cell.

\begin{table}[ht]
\centering
\small
\begin{tabular}{@{}l >{\raggedright\arraybackslash}p{0.42\linewidth} l l@{}}
\toprule
rung & change & value & step (exact)\\
\midrule
start & first run on the $7/8$ path with the menu of the earlier route-P stack, offset $2\cdot10^{-6}$ & $0.117708313410954$ & \\
E1 & start offset $2\cdot10^{-6}\to10^{-7}$ & $0.117708137107302$ & $-44075913/(2.5\cdot10^{14})$\\
E2 & existing height tables ($0.97$, $0.94$, $0.91$ times the plan, and a per-cell plan with multipliers $0.928$--$1.067$ at its knots) with $r=3/5$, $\lambda<2$ kernels & $0.117312236761635$ & $-395900345667/10^{15}$\\
E3 & new uniform height sets at $0.88$ (cells $0$--$140$) and $0.91$ (cells $0$--$199$) times the plan & $0.116174796162911$ & $-284360149681/(2.5\cdot10^{14})$\\
E4 & per-cell heights at the producer's moment-cap limit ($0.860$--$0.903$ times the plan, cells $0$--$250$) & $0.116110902986179$ & $-15973294183/(2.5\cdot10^{14})$\\
F1 & $79$ new $r=3/5$ kernels on the floors of (M1) & $0.116091900099009$ & $-1900288717/10^{14}$\\
F2 & the same kernels on the floors of (M2) & $0.116091488672439$ & $-41142657/10^{14}$\\
\midrule
total & & & $-323364947703/(2\cdot10^{14})$\\
\bottomrule
\end{tabular}
\caption{Re-tuning of the run of Section~\ref{sec:run2}. Each value is the exact output of one
run; all values depend on (H1)--(H4).}
\label{tab:ladder}
\end{table}

\subsection{Checks}\label{sec:checks}

For value (ii) the consumer checks, on every substep: admissibility $v\le\rho^2$ and
$hr<2$; the cap by Raymond's rule; $c>0$ and $0<c\delta<1$; the stored step with the
degree-$18$ bound and upward rounding; no crossing of $\qe$ before the terminal substep;
heights in $[\frac12,\frac52]$; $q_0>9/16$ and (S1); the terminal enclosure; and agreement
of two option generators. All report $0$ failures. Further checks:
\begin{itemize}[leftmargin=1.6em]
\item a step verifier and a from-scratch verifier that shares no code with the consumer and
re-derives every floor from the raw table files: both $0$ failures; the from-scratch
verifier was re-run on the archived per-substep dump for this paper ($38$\,s) with the same
result;
\item a second certification of the kernels by Arb with adaptive bisection, sharing no code
with the Bernstein certifier: $79/79$ new kernels at $64$ bits ($549302$ boxes), all $59$ rows
used by the run at $128$ bits ($425359$ boxes);
\item a re-certification of all $289$ (set, cell) pairs of the floors of (M1) at $224/320$-bit
lanes and different subdivisions: every one of the $3179$ mode floors lies $1.9\cdot10^{-4}$ to
$4.2\cdot10^{-4}$ above the value used, none below;
\item a recheck in $256$-bit Arb of Raymond's per-row source premises at every (cell, height,
source) used: $468$ triples, $0$ failures, minimum $k=t_L\log N_0/2-(1-h)/2$ equal to $0.43918$ ($t_L$ the left end of the cell),
minimum $h\log N_0$ equal to $12.236$, largest error sum in AF~(5.4) equal to
$1.4994\cdot10^{-7}<10^{-6}$;
\item the Lean kernel replay of every substep (Section~\ref{sec:lean}).
\end{itemize}
For value (i), the run was re-executed for this paper ($34$\,s) and reproduced $U$, the
substep count and the terminal cell.

\subsection{Proof of Theorem~\ref{thm:main}}

(i) The run of Section~\ref{sec:run1} satisfies (C1)--(C6) of Theorem~\ref{thm:barrier} with
$\theta=7/8$: these are exact rational checks performed by the consumer, by the independent
replay and by the Lean kernel. Its floor entries and kernel rows are those of (H2), and (W)
is part of (H2), with (H3) used inside it. By (H1) and Theorem~\ref{thm:barrier},
$\LB\le T^*+\qe/2\le U$, since $U$ is computed from upper bounds for $e^{-x}$ and
$\log(1+z)$ and rounded up.

(ii) The same argument applies to the run of Section~\ref{sec:run2}, whose floor entries and
kernel rows are those of (H4). \qed

\subsection{Results}

\begin{table}[ht]
\centering
\small
\begin{tabular}{@{}>{\raggedright\arraybackslash}p{0.44\linewidth} l >{\raggedright\arraybackslash}p{0.27\linewidth}@{}}
\toprule
chain and start & value & conditions\\
\midrule
de~Bruijn, $7/8$ half-plane alone (Lemma~\ref{lem:subst}) & $9/32=0.28125$ & (H1)\\
Raymond's chain, classical start \cite{Raymond26} & $0.149395073827688$ & (H2), (H3)\\
route~P, classical start (reference run) & $0.148250146775075$ & (H2), (H3), earlier route-P certificates\\
Raymond's chain, $11/12$ start & $0.129826438865344$ & $11/12$ half-plane, (H2), (H3)\\
route~P, $11/12$ start, not re-tuned & $0.128829131217197$ & $11/12$ half-plane, (H2), (H3), earlier route-P certificates\\
\textbf{Raymond's chain, $7/8$ start: Theorem~\ref{thm:main}(i)} & $\mathbf{0.118505710922701}$ & (H1)--(H3)\\
route~P, $7/8$ start, first run & $0.117708313410954$ & (H1)--(H3), earlier route-P certificates\\
\textbf{route~P re-tuned, $7/8$ start: Theorem~\ref{thm:main}(ii)} & $\mathbf{0.116091488672439}$ & (H1)--(H4)\\
\midrule
accepted peer-reviewed bound \cite[Corollary~2]{PlattTrudgian21} & $0.2$ & \\
\bottomrule
\end{tabular}
\caption{Values at the Platt--Trudgian height and the hypotheses each uses. The $11/12$ half-plane follows
from the $7/8$ one and is also the subject of a separate write-up \cite{OpenAI26QRH1112}.}
\label{tab:values}
\end{table}

On Raymond's chain the start alone moves the value from $0.149395073827688$ to
$0.118505710922701$, a difference of $0.030889362904987$. The re-tuned menu lowers it by
$1207111125131/(5\cdot10^{14})=0.002414222250262$ more.

% ===============================================================
\section{The value as a function of \texorpdfstring{$\theta$}{theta}}\label{sec:theta}
% ===============================================================

Since the chain reads the half-plane only through $q_0=(2\theta-1)^2+\epsilon$, it can be run
for hypothetical abscissas $\theta$. Table~\ref{tab:theta} lists Raymond's chain at the
Platt--Trudgian height with $\epsilon=2\cdot10^{-6}$ and the re-tuned menu of
Section~\ref{sec:run2} with $\epsilon=10^{-7}$. Every $\theta$ other than $7/8$ is
hypothetical, and each entry is conditional on the half-plane at that $\theta$, on (H2) and
(H3), and for the last column on (H4)-type certificates at the heights used. An independent
$192$-bit integrator written by the internal referee reproduces all $14$ values of the third
column.

\begin{table}[ht]
\centering
\small
\begin{tabular}{@{}l l l r l l@{}}
\toprule
$\theta$ & $(2\theta-1)^2/2$ & Raymond's chain & end cell & AF II substeps & re-tuned run\\
\midrule
$7/8$ & $0.28125$ & $0.118505710922701$ & $474$ & $0$ & $0.116091488672439$\\
$43/50$ & $0.2592$ & $0.114167305308158$ & $456$ & $124$ & \\
$17/20$ & $0.245$ & $0.111177850880042$ & $444$ & $270$ & $0.108886356046385$\\
$21/25$ & $0.2312$ & $0.108109947418064$ & $432$ & $571$ & \\
$41/50$ & $0.2048$ & $0.101732525557584$ & $406$ & $1944$ & \\
$4/5$ & $0.18$ & $0.095022125299536$ & $380$ & $3078$ & $0.093069284680992$\\
$39/50$ & $0.1568$ & $0.087983897822374$ & $351$ & $2869$ & \\
$3/4$ & $0.125$ & $0.076825435581115$ & $307$ & $1372$ & \\
$7/10$ & $0.08$ & $0.056949250938398$ & $227$ & $0$ & \\
$13/20$ & $0.045$ & $0.036622545758074$ & $146$ & $0$ & \\
$3/5$ & $0.02$ & $0.018149345876410$ & $72$ & $0$ & \\
$11/20$ & $0.005$ & $0.004876352555412$ & $19$ & $0$ & \\
$13/25$ & $0.0008$ & $0.000797732865289$ & $3$ & $0$ & \\
$101/200$ & $0.00005$ & $0.000050986694974$ & $0$ & $0$ & \\
\bottomrule
\end{tabular}
\caption{Values at the Platt--Trudgian height against a hypothetical zero-free abscissa
$\theta$. Decimals are exact expansions of the computed rationals (listed in
\texttt{PARAMETERS.md}). ``AF II substeps'' counts the substeps that use Raymond's sharp
continuum law; disabling it raises the affected values by at most $2.9\cdot10^{-7}$. At
$\theta=101/200$ the chain exceeds the de~Bruijn value because of the start offset, and the
de~Bruijn value is the better statement.}
\label{tab:theta}
\end{table}

The following are measured on this grid. The value decreases strictly in $\theta$. Near
$7/8$ the slope is about $0.289$; it peaks at about $0.407$ on $[0.65,0.70]$. No constraint
of the chain binds for $\theta\le7/8$: the terminal cell falls with $\theta$, the value
stays below the wall horizon $41/250$, and kernel admissibility becomes easier as $q$
shrinks.

The share of the descent carried by the floor-dependent term, $1-2U/q_0$, is $0.579$ at
$7/8$, $0.472$ at $0.80$, $0.288$ at $0.70$ and $0.093$ at $0.60$. Below $\theta\approx0.6$ the
value is within $10\%$ of $(2\theta-1)^2/2$, which uses no verified height. By linear
interpolation, $\LB<0.10$ at this height would need $\theta\approx0.815$ on Raymond's chain
or $\theta\approx0.82$ with the re-tuned menu.

At $\theta=7/8$, Raymond's chain with a head at $10X$, $100X$ and $1000X$ (heights at which RH
has not been verified) gives $0.112235848057212$, $0.107056451546934$ and
$0.102615591280579$. At the Platt--Trudgian height, $\theta=0.8535$ gives
$0.112233029447535$, so near $7/8$ a tenfold verified height is worth about $0.0215$ of
$\theta$.

% ===============================================================
\section{What limits the exponent \texorpdfstring{$7/8$}{7/8}}\label{sec:limits}
% ===============================================================

This section summarises internal analyses of the $7/8$ manuscript \cite{OpenAI26QRH}
\cite{IntQrhRecords}. Each item was checked in adversarial internal referee passes, and none
has been read outside the project. None improves $7/8$, and Theorem~\ref{thm:main} does not
depend on any of them. Equation and lemma numbers refer to \cite{OpenAI26QRH}.

\begin{enumerate}[leftmargin=1.8em]
\item \emph{Exponent bookkeeping.} From the displayed formulas (eqs.~(10.3), (12.4), (15.6),
(15.9)), the abscissa is $\sigma_0=1-h/6+b/12=\frac56+\frac M{12}-\frac\ell6$ with
$h=1-l_x+\ell$, where $M=l_x+l_y$ and $b=l_y-l_x$. On $M+\ell=1$ this is
$\sigma_0=\frac{11}{12}-\frac\ell4$. The manuscript's geometry
$(l_x,l_y,\ell)=(17/48,23/48,1/6)$ gives $7/8$; $\ell=0$ gives the $11/12$ of the first stage.
\item \emph{Joint limit.} The coefficient $1/12$ of $b$ in $\sigma_0$ comes from the
sixth-power exceptional term $P_a^{1/6}$ of the additive Gram bound (Proposition~15.2). With
it, the zero-loss range $2M\ge1+4\ell$ of Lemma~15.1 gives
$\sigma_0\ge\frac56+\frac{1+4\ell}{24}-\frac\ell6=\frac78$, with equality where that range is
tight. The high side binds at the row bin of Lemma~20.2 at $x=\frac12$,
$\delta^*=\frac{185}{138}-\frac{35\sqrt{16354}}{4692}\approx0.38664$, with slack
$\frac{35\sqrt{16354}}{12696}-\frac{1491}{4232}\approx2.2815\cdot10^{-4}$ (the manuscript
certifies $49/440640$). Re-optimising the manuscript's own inequalities gives at most
$\sigma_0\approx0.874957$ (a model optimum, not a theorem).
\item \emph{Density-hypothesis counts.} With Proposition~15.2 as stated, the floor bin and the
low side meet at $(\ell,M,b)=(\frac15,\frac45,0)$, giving $\frac{13}{15}$ even with
density-hypothesis row counts.
\item \emph{The exceptional term.} Three restructurings of the Cauchy--Schwarz step of
Proposition~15.3 were examined: subtracting the exceptional component before
Cauchy--Schwarz, H\"older-type bounds, and restricting or re-weighting the rows. None lowers
the abscissa. At main-term level (error terms sketched only), the term is a positive dual main
term of the row mean square; the paper's column
coefficient equals its own exceptional component, so subtraction is vacuous in that form; any
bound that sees the dual vector only through Lemma~15.1 is at least $\|A\|_2V$; and row
restrictions that remove the term also remove the high-side signal.
\item \emph{Height.} Given (H3), a half-plane valid only above $3\cdot10^{12}$ is equivalent
to the global statement, and the chain reads only $\sup(2\beta-1)^2$ over
$\gamma>(X+1)/2$; a height-dependent improvement is not an easier target.
\end{enumerate}

% ===============================================================
\section{Lean status}\label{sec:lean}
% ===============================================================

Three Lean~4 \cite{Lean4} developments bear on the result. All use Lean~4.34.1 and Mathlib
\cite{Mathlib20} commit \texttt{d13f23b7}. Every theorem below depends only on the axioms
\texttt{propext}, \texttt{Classical.choice} and \texttt{Quot.sound}; there is no
\texttt{sorry}, no \texttt{native\_decide} and no \texttt{axiom} declaration.

\paragraph{The half-plane (H1).} The theorem
\texttt{OAI.riemannZeta\_ne\_zero\_of\_seven\_eighths\_lt\_re}, stating
\texttt{riemannZeta s} $\ne0$ for $7/8<\Rea s$, was built from the initial commit
\texttt{adc7f1241} of \cite{OpenAI26QRH}: $7061$ build jobs, all $2924$ modules of the import
closure, $0$ errors. Its type equals the repository's comparator challenge elaborated against
Mathlib alone, and \texttt{riemannZeta} is Mathlib's function. A source scan of the closure
($486490$ lines) found no \texttt{sorry}, \texttt{admit}, \texttt{axiom},
\texttt{native\_decide}, \texttt{unsafe}, \texttt{implemented\_by}, \texttt{extern} or
\texttt{set\_option}. The repository update of 8~October~2026 (commit \texttt{301488868},
merged as \texttt{fd4aeeb2}) changes none of the $2924$ modules of this closure, the comparator
file, the lakefile, the manifest, the toolchain or the family-003 manuscripts (file list of the
commit compared with the closure on 9~October~2026); the build was not repeated at the new
commit. Not done: the official comparator run, an independent kernel replay
(\texttt{lean4checker} or nanoda) of the closure, and refereeing of the manuscript.

\paragraph{Lemma~\ref{lem:subst}, Lemma CG, Lemma R.} The package \texttt{raymond-core}
proves Lemma~\ref{lem:subst} (\texttt{RaymondCore.Substitution.substitution}) and
$\LB\le(2\theta-1)^2/2$ for the concrete $H_t$ of LeanCert \cite{LeanCert}, which formalises
$H_t$, the identity $H_0(z)=\frac18\xi(\frac12+\frac{iz}2)$, de~Bruijn's strip contraction
and the threshold characterisation of $\LB$. At $\theta=7/8$ it proves $\LB\le9/32$ from the
openai/math theorem. It also proves the coefficient monotonicity (Lemma CG) used in
Raymond's main eqs.~(11)--(13), and the analytic core of the phase estimate (Lemma R) used
by the floors (M1) of Section~\ref{sec:run2}, with Raymond's Stirling bounds AF~(4.6) as
explicit hypotheses.

\paragraph{The theorem.} The package \texttt{qrh-lambda} proves
\begin{center}
\small
\begin{tabular}{@{}l@{}}
\texttt{theorem QrhLambda.lambda\_le\_of\_qrh\_raymond (h : RaymondChainHyps)}\\
\texttt{\ \ \ \ (f : FloorTableHyps) : deBruijnNewmanLambda $\le$ 116091488672439 / 10 \^{} 15}
\end{tabular}
\end{center}
and the corresponding theorem with $118505710922701/10^{15}$ and Raymond's tables. Here
\texttt{deBruijnNewmanLambda} is LeanCert's
$\inf\{t:\text{every zero of }H_t\text{ is real}\}$. \texttt{RaymondChainHyps} has three
fields: (W); first contact (a continuous positive barrier that starts above the tail zeros
and is not above them later has a first contact time with a zero on it); and contact
exclusion (Lemma~\ref{lem:raymond} together with Proposition~\ref{prop:dyn}(ii), stated for
one substep). \texttt{FloorTableHyps} asserts (J) for each of the $468$ floor entries and (K)
for each of the $59$ kernel rows. Lean proves the start condition from the openai/math
theorem via Lemma~\ref{lem:subst}, the strict entry, the substep barrier lemma from the two
contact hypotheses, the terminal step, the exponential and logarithm bounds, the cap lemma,
and the soundness of a checker that recomputes each run from its start. Every substep of
both runs ($59439$ and $60675$) is replayed by \texttt{decide +kernel} in $30$ and $31$ chunk
theorems. Negative controls (for example $U$ lowered by $10^{-15}$, or one floor entry
lowered by $10^{-12}$) are rejected by the kernel.

An internal referee rebuilt the package from scratch, replayed the modules with
\texttt{leanchecker}, re-implemented the checker in Python (with $0$ mismatches against the
run over all $59439$ substeps), and proved that the hypotheses follow from RH, so they are
consistent unless RH is false.

\paragraph{What is not in Lean.} The three fields of \texttt{RaymondChainHyps}, including
(H3); (J) for each floor entry and (K) for each kernel row; and (H1) beyond the local build.
The Lean theorem therefore depends on the same inputs as Theorem~\ref{thm:main}. The
packages depend on machine-local builds and are not reproducible from a clean clone without
them.

% ===============================================================
\section{Dependencies and limitations}\label{sec:limitations}
% ===============================================================

\begin{table}[ht]
\centering
\small
\begin{tabular}{@{}>{\raggedright\arraybackslash}p{0.2\linewidth}>{\raggedright\arraybackslash}p{0.47\linewidth}>{\raggedright\arraybackslash}p{0.25\linewidth}@{}}
\toprule
condition & status & without it\\
\midrule
(H1) $\zeta\ne0$ on $\Rea s>7/8$ & claim of \cite{OpenAI26QRH}; Lean proof rebuilt locally, standard axioms; comparator and independent kernel replay not run; not peer reviewed & the chain gives $0.149395073827688$\\
(H2) Raymond's chain & unrefereed preprint \cite{Raymond26}; two internal review passes \cite{IntReviewA,IntReviewB}, no false step found, expository gaps with written fixes; no reader outside the project & (H1) alone gives $9/32$\\
(H3) RH to $3000175332800$ & peer reviewed \cite{PlattTrudgian21}; used once, at $t=0$ & no head wall\\
(H4) route-P certificates & project-internal, mostly single implementation, one internal referee pass & (ii) reduces to (i)\\
\bottomrule
\end{tabular}
\caption{The four conditions.}
\label{tab:deps}
\end{table}

The internal reviews of (H2) found expository gaps, each with a written fix: the counting
step of the head localisation (AF~\S2); the ideal model and error budget of the wall, which
are only in code (AF~\S3); the endpoint $t=0$ of the wall; the uniformity of the floors in
$x$ (AF~\S5); the derivation of main eq.~(1); the start written only for $q_0>1$; and the
compressed first-contact argument, written out here as Theorem~\ref{thm:barrier}.

The weakest items in (H4) are the structured floors (M3) of Section~\ref{sec:run2}, used on
$11304$ substeps. They rest on an outer-lemma inequality certified entry by entry by
project-written engines (single implementations with internal referee passes), on a measured, uncertified allowance of $0.0023$ in the time-box step, on a lemma for
uniformity in $x$ that has not been reviewed, and on a heuristic allowance
$10^{-12}\cdot\Sigma D\cdot\Sigma\mu$ per entry for the floating-point FFT used in their
tails, in place of a proved error bound. The joint-pair and gain code of (M1)--(M2) is a single implementation.
The re-planned heights rely on a licence to combine floors from several height tables in one
run, which has not been refereed outside the project.

The run scripts, tables and dumps are in an untracked working tree; the package accompanying
this draft pins them by SHA-256. None of the four conditions is established by this paper, and
Theorem~\ref{thm:main} adds no bound on $\LB$ to the literature until they are.

\section*{Acknowledgements}

The barrier method, its floor estimates, kernel catalogue and comparison arithmetic are
D.~J.~Raymond's \cite{Raymond26}; this paper changes the start of his construction and adds
certificates to it. His deposit, which includes the TeX source, the analytic companion and a
full replay archive, made the internal reviews and the computations possible. The half-plane
$\Rea s>7/8$ and its Lean proof are from the openai/math repository \cite{OpenAI26QRH}. The
verified height is Platt and Trudgian's \cite{PlattTrudgian21}, and the zero dynamics,
confinement and normalisation are from the Polymath project \cite{Polymath19}. The Lean
developments build on Mathlib \cite{Mathlib20} and LeanCert \cite{LeanCert}; the interval
computations use Arb and FLINT \cite{JohanssonArb17,FLINT}.

\appendix
\section{Reproduction}\label{app:repro}

Paths are relative to the project repository. The scripts and run artefacts are in its
untracked \texttt{tmp/} tree; their SHA-256 hashes are listed in \texttt{PROVENANCE.md} of
the accompanying package. Expected outputs are given after each block.

\paragraph{Value (i)} (about $35$\,s).
{\small
\begin{verbatim}
.venv-arb/bin/python tmp/oai-lambda-raymond.py OUT.json X1 --thetas 7/8
\end{verbatim}}
\noindent Expected: \texttt{U = 118505710922701/1000000000000000}, $60675$ substeps, last
row $474$. The internal referee's independent implementation is
\texttt{tmp/oai-lambda-ref-raymond.py}.

\paragraph{Value (ii)} (consumer about $11$\,min, each verifier about $1$\,min).
{\small
\begin{verbatim}
cd tmp
XK=qrh-tune-out/xkernels_final.json KK=qrh-tune-out/xkernels_final.json \
  ./qrh-tune-final.sh F2_repro 1/10000000 \
  W4 m097 m094 m091 QA_m0.88 QB_m0.91 QC_cliff
./qrh-ref-run.sh refverify /tmp/qrh-tune-dump-F2_repro F2_repro_refverify
\end{verbatim}}
\noindent Expected: \texttt{U = U\_robust = 116091488672439/1000000000000000}, $59439$
substeps, terminal cell $464$, $0$ failures in both the consumer's step verifier and the
from-scratch verifier.

\paragraph{The $\Lambda(\theta)$ table.}
{\small
\begin{verbatim}
tmp/qrh-theta-curve-raymond.sh
tmp/qrh-theta-curve-routeP.sh      # output: tmp/qrh-theta-curve-out/table.json
\end{verbatim}}

\paragraph{Lean.} With the dependency builds described in \texttt{LEAN.md}:
{\small
\begin{verbatim}
cd formalization/lean/qrh-lambda
scripts/build_chunks.sh F2 3; scripts/build_chunks.sh R78 3
lake build
lake env lean scripts/AxiomsCheck.lean
scripts/check-oai.sh
lake env lean scripts/NegativeControl.lean
\end{verbatim}}

\bibliographystyle{amsplain}
\bibliography{refs}

\end{document}
